\documentclass[11pt,letterpaper]{article}

% ---------------------------------------------------------------
%  Packages
% ---------------------------------------------------------------
\usepackage[margin=1in]{geometry}
\usepackage{amsmath,amssymb,amsthm,mathtools}
\usepackage{bm}
\usepackage{enumitem}
\usepackage{hyperref}
\usepackage{xcolor}
\usepackage{tikz}
\usepackage{booktabs}
\usepackage{microtype}
\usepackage{titlesec}
\usepackage{parskip}

\hypersetup{
    colorlinks=true,
    linkcolor=blue!50!black,
    citecolor=blue!50!black,
    urlcolor=blue!60!black
}

% ---------------------------------------------------------------
%  Theorem environments
% ---------------------------------------------------------------
\newtheorem{theorem}{Theorem}[section]
\newtheorem{lemma}[theorem]{Lemma}
\newtheorem{proposition}[theorem]{Proposition}
\newtheorem{corollary}[theorem]{Corollary}
\theoremstyle{definition}
\newtheorem{definition}[theorem]{Definition}
\newtheorem{example}[theorem]{Example}
\newtheorem{remark}[theorem]{Remark}
\newtheorem*{notation}{Notation}

% ---------------------------------------------------------------
%  Macros
% ---------------------------------------------------------------
\newcommand{\C}{\mathbb{C}}
\newcommand{\R}{\mathbb{R}}
\newcommand{\N}{\mathbb{N}}
\newcommand{\Z}{\mathbb{Z}}
\newcommand{\F}{\mathbb{F}}
\newcommand{\Sphere}{\mathbb{S}}
\newcommand{\Polygram}{\mathcal{P}}
\newcommand{\Mset}{\mathcal{M}}
\newcommand{\Hilb}{\mathcal{H}}
\newcommand{\Cov}{\mathrm{N}}
\newcommand{\vol}{\operatorname{vol}}
\newcommand{\bdry}{\partial}
\newcommand{\diam}{\operatorname{diam}}
\newcommand{\dist}{\operatorname{dist}}
\newcommand{\range}{\operatorname{range}}
\newcommand{\im}{\operatorname{Im}}
\newcommand{\re}{\operatorname{Re}}
\newcommand{\rank}{\operatorname{rank}}
\newcommand{\tr}{\operatorname{tr}}
\newcommand{\inner}[2]{\langle #1 \mid #2 \rangle}
\newcommand{\ket}[1]{\lvert #1 \rangle}
\newcommand{\bra}[1]{\langle #1 \rvert}
\newcommand{\norm}[1]{\lVert #1 \rVert}
\newcommand{\abs}[1]{\lvert #1 \rvert}
\newcommand{\Set}[1]{\left\{#1\right\}}
\DeclareMathOperator*{\argmin}{arg\,min}
\DeclareMathOperator*{\argmax}{arg\,max}

\titleformat{\section}{\Large\bfseries}{\thesection}{0.7em}{}
\titleformat{\subsection}{\large\bfseries}{\thesubsection}{0.6em}{}

% ---------------------------------------------------------------
%  Title
% ---------------------------------------------------------------
\title{\textbf{Polygrams:\\
A Theoretical Treatment of Kronecker-Succinct Vector Families}}
\author{}
\date{May 2026}

\begin{document}
\maketitle

\begin{abstract}
We introduce and study \emph{Polygrams}, a family of implicit unit-vector
representations indexed by tuples of points on the complex projective line
$\C P^1$ whose expansion is a Kronecker product over $\C^2$. A Polygram of
\emph{order} $n = M + k$ stores only $2n$ real parameters yet represents a
distinguished element of the unit sphere of $\C^{2^{n}}$, an ambient space whose
dimension is exponential in $n$. The pay-off is that inner products,
similarities, and a wide range of geometric queries factor across the parameter
tuple and therefore cost $\Theta(n)$ rather than $\Theta(2^{n})$. We develop the
basic theory of these objects from first principles in a tutorial style. After
establishing norm preservation and the factorised-inner-product identity, we
study (i) the geometry of the Polygram manifold inside the unit sphere, with
explicit dimension and codimension counts; (ii) covering-number based
\emph{approximation lower bounds} that quantify how badly arbitrary states can
be approximated by Polygrams; (iii) Lipschitz \emph{stability} of the expansion
map under parameter perturbation, including condition numbers for similarity
queries; (iv) the \emph{gauge group} acting on the parameter space and a
constructive identifiability result requiring only $4n + 1$ overlap measurements;
and (v) a formal comparison to matrix product states (MPS) and tensor trains
(TT), showing that Polygrams are precisely the bond-dimension-one case and
analysing the resulting expressivity-versus-storage trade-off. We close with
algorithms for fitting and merging Polygrams (coordinate descent, Riemannian
gradient flow) and a list of open problems, including refined covering
bounds, sample complexity for parameter recovery from noisy overlaps, and
sharper expressivity separations from MPS at small bond dimension.
\end{abstract}

\tableofcontents
\bigskip
\hrule
\bigskip

% ===============================================================
\section{Introduction}
% ===============================================================

Many problems in machine learning, similarity search, signal processing, and
quantum-inspired classical computation are organised around \emph{collections of
unit vectors in a very high-dimensional space}. Word embeddings, locality-
sensitive sketches, kernel feature maps, randomised projections, dictionary
atoms, content-addressable hashes, and quantum state vectors are all elements
of $\Sphere(\Hilb)$ for some Hilbert space $\Hilb$, and the operations we
usually care about, namely inner products, projections, and nearest-neighbour
queries, are linear or quadratic in $\Hilb$.

When $\dim\Hilb$ is modest there is nothing more to say: write the vector down
as an array of $\dim\Hilb$ numbers and use BLAS. When $\dim\Hilb$ is huge the
situation is different. Storing $\Hilb$ explicitly is impossible even for
$\dim\Hilb = 2^{40}$, and yet we routinely \emph{simulate} vectors in such
spaces by exploiting structure. Sparse vectors, low-rank matrices, low-degree
polynomials, random sketches, and tensor networks are all examples of
\emph{succinct} families: each is a parameterised subset of $\Hilb$ whose
elements admit asymptotically optimal storage and whose key operations
(addition, inner product, action of a linear operator) factor through the
parameterisation.

This paper introduces a particular succinct family, the \emph{Polygram}, whose
defining property is that its elements are pure tensor products of qubit-style
vectors in $\C^{2}$:
\[
\bm{f} \;=\; \bm{v}_{1} \otimes \bm{v}_{2} \otimes \cdots \otimes \bm{v}_{n},
\qquad \bm{v}_{i} \in \Sphere(\C^{2}).
\]
A Polygram of order $n$ stores $n$ such factors and therefore $2n$ real
parameters (each $\bm{v}_{i}$ is a point on the Bloch sphere $\C P^{1}\cong
\Sphere^{2}$), and yet the vector it represents is an element of the unit sphere
of $\C^{2^{n}}$. The trick is venerable. It appears in physics under the name
\emph{product state} or \emph{Hartree ansatz}, in tensor-network theory as a
matrix product state with bond dimension one, in algebraic statistics as a
rank-one element of $(\C^{2})^{\otimes n}$, and informally throughout numerical
analysis as a tensor outer product. What this paper offers is not the trick,
which is over a century old, but a tutorial-style theoretical study of its
properties when we treat the parameter tuple as a first-class data structure
intended for algorithmic use.

% ---------------------------------------------------------------
\subsection{Why a new name?}
% ---------------------------------------------------------------

A reader from quantum information theory may protest that ``Polygram'' is just
a new name for a separable pure state on $n$ qubits. That is correct, but the
intended use is different and the change of name signals that.

Treating product states as a generic engineering primitive forces a particular
algorithmic discipline. We refuse to materialise the tensor, we benchmark our
algorithms against the cost of similarity queries on explicit
$2^{n}$-dimensional arrays, we worry about identifiability under noisy
overlaps, and we measure expressivity by the worst-case distance from an
arbitrary unit vector to the Polygram manifold. None of these questions is
novel in absolute terms, but combining them under a single succinct-vector
abstraction is useful, in the same way that ``hash table'' is useful even
though every hash table is just an array equipped with a function.

We adopt the term \emph{Polygram} because it suggests an object built out of
many small pieces (``poly-'') that together write something (``-gram''). The
written object is the parameter tuple; the thing it writes is the implicit
vector $\bm{f}$.

% ---------------------------------------------------------------
\subsection{Informal preview of the main results}
% ---------------------------------------------------------------

Throughout the paper $n = M + k$ denotes the total number of factors. The
distinction between $M$ \emph{base} factors and $k$ \emph{amplitude} factors
will reappear in algorithmic sections but is irrelevant to the mathematical
content of the first results: a Polygram is determined by its $n$ factors and
nothing else. We use $D = 2^{n}$ to denote the ambient dimension and we work
exclusively over $\C$.

\begin{itemize}[leftmargin=*,itemsep=4pt]
\item \textbf{Norm and inner product.} Every Polygram has unit norm, and the
      inner product between two Polygrams factors as
      $\inner{\bm{f}}{\bm{g}} = \prod_{i=1}^{n} \inner{\bm{v}_{i}}{\bm{w}_{i}}$.
      Both facts follow from elementary properties of the Kronecker product,
      but together they imply that all similarity queries cost $\Theta(n)$ time
      and $\Theta(n)$ space, an exponential improvement over the explicit
      representation.
\item \textbf{Geometry.} The set of order-$n$ Polygrams, viewed modulo the
      $U(1)$ global phase, is a real $2n$-dimensional smooth submanifold of
      $\C P^{D-1}$. The ambient projective space has real dimension $2D - 2 =
      2^{n+1} - 2$, so for $n \geq 2$ the Polygram manifold has codimension
      $2^{n+1} - 2n - 2$, which grows exponentially in $n$. Polygrams are a
      vanishingly thin sliver of unit-vector space.
\item \textbf{Approximation.} The fact that Polygrams are a small submanifold
      forces a sharp negative result. We prove that for every $n \geq 2$ there
      exist unit vectors in $\C^{2^{n}}$ whose best Polygram approximation has
      squared overlap at most $C \cdot n \cdot 2^{-n}$ for some absolute
      constant $C$. The argument is a volume/covering computation; the
      conclusion is that Polygrams are \emph{not} a universal approximator and
      should be combined with other primitives (e.g.\ low-rank mixtures) when
      universal approximation matters.
\item \textbf{Stability.} The expansion map from parameters to vectors is
      $\sqrt{n}$-Lipschitz with respect to the natural Euclidean metric on the
      parameter torus and the Euclidean metric on $\C^{D}$. Similarity queries
      are $1$-Lipschitz, so small parameter perturbations produce only small
      changes in similarity. The factorised structure makes the Jacobian and
      Hessian of any per-pair query computable in $\Theta(n)$ time.
\item \textbf{Identifiability.} The continuous redundancy in the parameter
      tuple is exactly the diagonal $U(1)$ acting by global phase, plus a
      residual $U(1)^{n-1}$ that rescales factors against one another. Modulo
      this gauge, the Polygram is determined by $4n + 1$ generic overlap
      measurements with chosen probe Polygrams, a result we prove
      constructively in Section~\ref{sec:identifiability}.
\item \textbf{Comparison to MPS/TT.} Polygrams are exactly matrix product
      states with bond dimension $D_{\mathrm{bond}} = 1$. Increasing the bond
      dimension to $D_{\mathrm{bond}}$ increases storage to $O(n \cdot
      D_{\mathrm{bond}}^{2})$ and similarity-query cost to $O(n \cdot
      D_{\mathrm{bond}}^{3})$ but yields strictly more expressivity. We make
      this trade-off precise in Section~\ref{sec:mps}.
\end{itemize}

% ---------------------------------------------------------------
\subsection{Reading guide}
% ---------------------------------------------------------------

The paper is intended to be read linearly. Sections \ref{sec:prelim}--
\ref{sec:basic} establish notation and prove the algebraic backbone of the
theory; nothing later depends on conventions introduced after
Section~\ref{sec:basic}. Section~\ref{sec:geometry} introduces the differential
geometry needed for the covering bound of Section~\ref{sec:approximation} and
the stability result of Section~\ref{sec:stability}; readers comfortable with
smooth manifolds can skim Section~\ref{sec:geometry}.
Section~\ref{sec:identifiability} is self-contained modulo the gauge
calculation. Sections \ref{sec:mps}--\ref{sec:algorithms} are mostly
independent of the theoretical core and can be read in either order. We close
with worked examples (Section~\ref{sec:examples}) and a list of open problems
(Section~\ref{sec:open}).

A reader purely interested in the engineering claim ``Polygrams give $\Theta(n)$
similarity queries on objects of effective dimension $2^{n}$, and here is what
that costs you in expressivity'' can read Sections~\ref{sec:basic},
\ref{sec:approximation}, and \ref{sec:mps} in that order.

% ===============================================================
\section{Preliminaries}
\label{sec:prelim}
% ===============================================================

We collect the conventions and elementary facts used throughout. None of this
material is original.

% ---------------------------------------------------------------
\subsection{Hilbert spaces and inner products}
% ---------------------------------------------------------------

We work with finite-dimensional complex Hilbert spaces. By default
$\Hilb \cong \C^{N}$ is equipped with the standard Hermitian inner product
\[
\inner{\bm{u}}{\bm{v}} \;=\; \sum_{i=1}^{N} \overline{u_{i}}\, v_{i},
\]
which is conjugate-linear in the first slot and linear in the second. The
Euclidean norm is $\norm{\bm{v}}^{2} = \inner{\bm{v}}{\bm{v}}$. The unit sphere
is $\Sphere(\Hilb) = \{\bm{v} \in \Hilb : \norm{\bm{v}} = 1\}$.

Two unit vectors $\bm{u}, \bm{v}$ are equivalent up to global phase, written
$\bm{u} \sim \bm{v}$, iff $\bm{u} = e^{i\alpha} \bm{v}$ for some $\alpha \in
\R$. The quotient $\Sphere(\Hilb) / {\sim}$ is the complex projective space
$\C P^{N - 1}$, with real dimension $2N - 2$. We will pass freely between
the unit sphere and projective space when the global phase is immaterial.

% ---------------------------------------------------------------
\subsection{Kronecker product}
% ---------------------------------------------------------------

Given vectors $\bm{u} \in \C^{p}$ and $\bm{v} \in \C^{q}$, the Kronecker (or
tensor) product $\bm{u} \otimes \bm{v} \in \C^{pq}$ is the vector indexed by
pairs $(i,j) \in [p] \times [q]$ with entries
\[
(\bm{u} \otimes \bm{v})_{(i,j)} \;=\; u_{i}\, v_{j}.
\]
We use lexicographic ordering of index pairs throughout. The Kronecker product
is bilinear and associative but \emph{not} commutative; the order of factors
is part of the data.

\begin{lemma}[Multiplicativity of the inner product under tensoring]
\label{lem:ip-tensor}
For $\bm{u}_{1}, \bm{u}_{2} \in \C^{p}$ and $\bm{v}_{1}, \bm{v}_{2} \in \C^{q}$,
\[
\inner{\bm{u}_{1} \otimes \bm{v}_{1}}{\bm{u}_{2} \otimes \bm{v}_{2}}
   \;=\; \inner{\bm{u}_{1}}{\bm{u}_{2}}\,\inner{\bm{v}_{1}}{\bm{v}_{2}}.
\]
\end{lemma}
\begin{proof}
Expand the inner product on the left:
\[
\sum_{i,j} \overline{(\bm{u}_{1} \otimes \bm{v}_{1})_{(i,j)}}
            (\bm{u}_{2} \otimes \bm{v}_{2})_{(i,j)}
   = \sum_{i,j} \overline{u_{1,i}}\,\overline{v_{1,j}}\,u_{2,i}\,v_{2,j}.
\]
The double sum factors as $\big(\sum_{i}\overline{u_{1,i}} u_{2,i}\big)
\big(\sum_{j}\overline{v_{1,j}} v_{2,j}\big) = \inner{\bm{u}_{1}}{\bm{u}_{2}}
\inner{\bm{v}_{1}}{\bm{v}_{2}}$, which is the right-hand side. \qed
\end{proof}

\begin{corollary}[Norm multiplicativity]
\label{cor:norm-mult}
$\norm{\bm{u} \otimes \bm{v}} = \norm{\bm{u}} \cdot \norm{\bm{v}}$.
\end{corollary}

By induction Lemma~\ref{lem:ip-tensor} extends to any number of factors. We
will use the $n$-fold version freely.

% ---------------------------------------------------------------
\subsection{The Bloch sphere}
% ---------------------------------------------------------------

A unit vector $\bm{v} \in \Sphere(\C^{2})$ can be written, up to global phase,
in the Bloch form
\[
\bm{v}(\theta, \phi) \;=\; \begin{pmatrix} \cos\theta \\
                                            e^{i\phi}\sin\theta \end{pmatrix},
\qquad \theta \in [0, \pi/2], \;\; \phi \in [0, 2\pi).
\]
Every point of $\C P^{1}$ has a unique pre-image in this parameterisation
except the poles $\theta \in \{0, \pi/2\}$, at which $\phi$ is ambiguous. The
parameterisation realises $\C P^{1}$ as the $2$-sphere $\Sphere^{2}$; identifying
$(\theta, \phi)$ with spherical coordinates gives the usual Bloch ball picture
familiar from quantum mechanics.

\begin{lemma}[Bloch inner product]
\label{lem:bloch-ip}
For $\bm{v}_{1} = \bm{v}(\theta_{1}, \phi_{1})$ and $\bm{v}_{2} = \bm{v}(
\theta_{2}, \phi_{2})$,
\[
\inner{\bm{v}_{1}}{\bm{v}_{2}}
   \;=\; \cos\theta_{1}\cos\theta_{2}
       + e^{i(\phi_{2} - \phi_{1})}\sin\theta_{1}\sin\theta_{2}.
\]
\end{lemma}
\begin{proof}
Compute coordinate-by-coordinate:
\[
\inner{\bm{v}_{1}}{\bm{v}_{2}}
   = \overline{\cos\theta_{1}} \cdot \cos\theta_{2}
   + \overline{e^{i\phi_{1}}\sin\theta_{1}} \cdot e^{i\phi_{2}}\sin\theta_{2}
   = \cos\theta_{1}\cos\theta_{2}
   + e^{i(\phi_{2} - \phi_{1})} \sin\theta_{1}\sin\theta_{2}.
\]
The $\theta_{i}$ are real, so the conjugation acts only on the $\phi$ phase.
\qed
\end{proof}

\begin{remark}
Lemma~\ref{lem:bloch-ip} immediately recovers the familiar identity
$\abs{\inner{\bm{v}_{1}}{\bm{v}_{2}}}^{2} = \cos^{2}\!\big(\Theta/2\big)$ where
$\Theta$ is the great-circle distance between the two corresponding points on
the Bloch sphere $\Sphere^{2}$.
\end{remark}

% ---------------------------------------------------------------
\subsection{Big-O conventions}
% ---------------------------------------------------------------

We use $f = O(g)$, $f = \Omega(g)$, $f = \Theta(g)$ in their standard
asymptotic meanings. ``Polylog'' means $O\big((\log N)^{c}\big)$ for some
constant $c$. ``Generic'' means ``true on a Zariski-open subset of full
measure''; we use this in identifiability and dimension statements. Absolute
constants are denoted $C, C', c, c', \dots$ and their values may change line to
line.

% ===============================================================
\section{Polygrams: definition and basic properties}
\label{sec:basic}
% ===============================================================

% ---------------------------------------------------------------
\subsection{Definition}
% ---------------------------------------------------------------

\begin{definition}[Polygram]
\label{def:polygram}
Let $n \geq 1$. An \emph{order-$n$ Polygram} is an ordered $n$-tuple
\[
\Polygram \;=\; (\bm{v}_{1}, \bm{v}_{2}, \dots, \bm{v}_{n})
\qquad \text{with} \qquad \bm{v}_{i} \in \Sphere(\C^{2}).
\]
Its \emph{expansion} is the unit vector
\[
\bm{f}(\Polygram) \;=\; \bm{v}_{1} \otimes \bm{v}_{2} \otimes \cdots \otimes
\bm{v}_{n} \;\in\; \Sphere(\C^{2^{n}}).
\]
Each factor is parameterised by Bloch angles $(\theta_{i}, \phi_{i})$ so the
parameter space is $\Mset_{n} = ([0, \pi/2] \times [0, 2\pi))^{n}$. The map
$\bm{f} : \Mset_{n} \to \Sphere(\C^{2^{n}})$ is called the \emph{expansion
map}; its image $\Sigma_{n} = \bm{f}(\Mset_{n})$ is the \emph{Polygram
variety}. We refer to $n$ as the \emph{order} of the Polygram and to $D = 2^{n}$
as its \emph{ambient dimension}.
\end{definition}

\begin{notation}
We retain the $M + k$ split from the introduction when it is operationally
relevant: $M$ ``base'' factors $(\theta_{i}^{b}, \phi_{i}^{b})_{i=1}^{M}$ and
$k$ ``amplitude'' factors $(\theta_{j}^{a}, \phi_{j}^{a})_{j=1}^{k}$, with
$n = M + k$. The split affects which factors are updated in particular
algorithms (e.g.\ dictionary learning may freeze the base factors while
optimising the amplitude factors), but no theorem in this paper uses the
distinction.
\end{notation}

% ---------------------------------------------------------------
\subsection{Norm preservation}
% ---------------------------------------------------------------

\begin{theorem}[Norm preservation]
\label{thm:norm}
For every Polygram $\Polygram$, the expansion $\bm{f}(\Polygram)$ has unit
Euclidean norm.
\end{theorem}
\begin{proof}
By Corollary~\ref{cor:norm-mult} applied iteratively,
\[
\norm{\bm{f}(\Polygram)}
   \;=\; \prod_{i=1}^{n} \norm{\bm{v}_{i}}
   \;=\; \prod_{i=1}^{n} 1
   \;=\; 1. \qed
\]
\end{proof}

This is essentially trivial, but it has a useful consequence: the expansion
map sends parameters directly into the unit sphere, with no normalisation
needed. In contrast, neural-network factorisations of large vectors typically
require a final $L^{2}$ normalisation that obscures derivatives.

% ---------------------------------------------------------------
\subsection{Factorised inner product}
% ---------------------------------------------------------------

This is the result that makes the entire data structure work.

\begin{theorem}[Factorised inner product]
\label{thm:fip}
Let $\Polygram = (\bm{v}_{1}, \dots, \bm{v}_{n})$ and
$\Polygram' = (\bm{w}_{1}, \dots, \bm{w}_{n})$ be order-$n$ Polygrams. Then
\[
\inner{\bm{f}(\Polygram)}{\bm{f}(\Polygram')}
   \;=\; \prod_{i=1}^{n} \inner{\bm{v}_{i}}{\bm{w}_{i}}.
\]
\end{theorem}
\begin{proof}
Apply Lemma~\ref{lem:ip-tensor} to $\bm{f}(\Polygram) = \bm{v}_{1} \otimes
\bm{v}_{2:n}$ and $\bm{f}(\Polygram') = \bm{w}_{1} \otimes \bm{w}_{2:n}$,
where $\bm{v}_{2:n}$ denotes the tensor product of factors $2, \dots, n$.
This yields
\[
\inner{\bm{f}(\Polygram)}{\bm{f}(\Polygram')}
   \;=\; \inner{\bm{v}_{1}}{\bm{w}_{1}} \cdot
        \inner{\bm{v}_{2:n}}{\bm{w}_{2:n}}.
\]
Induct on $n$. The base case $n = 1$ is immediate. The inductive step gives the
claim. \qed
\end{proof}

\begin{remark}[Concrete formula]
\label{rem:concrete-ip}
Combining Theorem~\ref{thm:fip} with Lemma~\ref{lem:bloch-ip} yields a closed
form in Bloch angles:
\[
\inner{\bm{f}(\Polygram)}{\bm{f}(\Polygram')}
   \;=\; \prod_{i=1}^{n} \bigg[
        \cos\theta_{i}\cos\theta'_{i}
      + e^{i(\phi'_{i} - \phi_{i})} \sin\theta_{i}\sin\theta'_{i}
   \bigg].
\]
\end{remark}

% ---------------------------------------------------------------
\subsection{Computational consequences}
% ---------------------------------------------------------------

\begin{corollary}[Similarity complexity]
\label{cor:complexity}
Let $\Polygram, \Polygram'$ be order-$n$ Polygrams stored as parameter tuples,
each occupying $O(n)$ machine words. Then:
\begin{enumerate}[label=(\roman*),leftmargin=*]
\item $\inner{\bm{f}(\Polygram)}{\bm{f}(\Polygram')}$ can be computed in
      $\Theta(n)$ arithmetic operations and $\Theta(n)$ memory accesses.
\item $\norm{\bm{f}(\Polygram) - \bm{f}(\Polygram')}^{2} = 2 - 2\re\,
      \inner{\bm{f}(\Polygram)}{\bm{f}(\Polygram')}$ inherits the same
      complexity.
\item The squared-overlap similarity
      $\abs{\inner{\bm{f}(\Polygram)}{\bm{f}(\Polygram')}}^{2}$ likewise has
      complexity $\Theta(n)$.
\end{enumerate}
The corresponding cost on explicit vectors of dimension $D = 2^{n}$ is
$\Theta(D) = \Theta(2^{n})$, an exponential improvement.
\end{corollary}
\begin{proof}
For (i), use the formula of Remark~\ref{rem:concrete-ip}: each of the $n$
factors costs $O(1)$, and multiplying them together costs $O(n)$. For (ii) and
(iii), apply (i) once.\qed
\end{proof}

\begin{corollary}[Gram-matrix complexity]
\label{cor:gram}
The $N \times N$ Gram matrix
$G_{ij} = \inner{\bm{f}(\Polygram^{(i)})}{\bm{f}(\Polygram^{(j)})}$ for $N$
order-$n$ Polygrams can be computed in $\Theta(N^{2} n)$ time and $\Theta(N^{2})$
storage. The explicit-vector cost is $\Theta(N^{2} \cdot 2^{n})$.
\end{corollary}

In particular, even when $D = 2^{n}$ exceeds machine memory, the Gram matrix
fits in memory and is computed in time independent of $D$. This is the main
practical pay-off of the data structure.

\begin{corollary}[Norm scaling under tensor products of Polygrams]
\label{cor:tensor-of-polygrams}
The tensor product $\bm{f}(\Polygram) \otimes \bm{f}(\Polygram')$ is itself a
Polygram of order $n + n'$ obtained by concatenating parameter tuples.
\end{corollary}
\begin{proof}
The Kronecker product of tensor products is the tensor product of the
concatenated tuples, by associativity. Norm and inner-product factorisation
both extend trivially. \qed
\end{proof}

The corollary implies that the class of Polygrams is closed under taking
products, which is occasionally convenient (Section~\ref{sec:algorithms}).

% ---------------------------------------------------------------
\subsection{Where the structure fails}
% ---------------------------------------------------------------

It is important to be honest about what Polygrams \emph{do not} do.

\begin{itemize}[leftmargin=*, itemsep=4pt]
\item \textbf{Closure under addition.} The set of Polygrams is not closed under
      vector addition. If $\bm{f}, \bm{g}$ are both Polygrams, then $\bm{f} +
      \bm{g}$ is generally \emph{not} a Polygram (it is a rank-$2$ tensor, not
      rank-$1$). Polygrams are a multiplicative, not additive, family.
\item \textbf{Closure under linear operators.} A general linear operator
      $A : \C^{D} \to \C^{D}$ applied to a Polygram does not preserve the
      Polygram structure. Only operators of the form $A = A_{1} \otimes \cdots
      \otimes A_{n}$ (local operators) preserve it. Operators outside this
      class force decompression.
\item \textbf{Sparsity.} The expansion $\bm{f}(\Polygram)$ is almost never
      sparse: typically every one of its $D$ coordinates is non-zero, and the
      coordinates are products of the per-factor coordinates.
\end{itemize}

These limitations are sharp. The right mental model is that Polygrams are a
\emph{multiplicative} compression scheme; they are powerful for queries that
factor multiplicatively and useless for queries that require sums.

% ===============================================================
\section{Geometry of the Polygram variety}
\label{sec:geometry}
% ===============================================================

We now study $\Sigma_{n} = \bm{f}(\Mset_{n})$ as a subset of the unit sphere
$\Sphere(\C^{D})$ where $D = 2^{n}$. The headline is that $\Sigma_{n}$ is a
smooth manifold of real dimension $2n + 1$ (counting global phase) inside an
ambient sphere of real dimension $2D - 1$. This dimension gap is the source of
all subsequent expressivity bounds.

% ---------------------------------------------------------------
\subsection{The Polygram manifold}
% ---------------------------------------------------------------

\begin{proposition}[Polygram manifold dimension]
\label{prop:dim}
The image $\Sigma_{n} = \bm{f}(\Mset_{n})$ is a real analytic submanifold of
$\Sphere(\C^{D})$ of real dimension $2n + 1$. The same set viewed modulo
global phase (i.e.\ $[\Sigma_{n}] \subset \C P^{D-1}$) is a real analytic
submanifold of real dimension $2n$.
\end{proposition}

\begin{proof}[Proof sketch]
The map $\bm{f}$ is real-analytic and its differential at a generic point has
full row rank. Specifically, in the projective quotient one can use the
parameterisation $\Sigma_{n} \cong (\C P^{1})^{n}$ given by
\[
\big([\bm{v}_{1}], \dots, [\bm{v}_{n}]\big) \;\longmapsto\;
[\bm{v}_{1} \otimes \cdots \otimes \bm{v}_{n}].
\]
The Segre embedding is a classical fact in algebraic geometry: $(\C P^{1})^{n}$
embeds as a smooth subvariety of $\C P^{2^{n} - 1}$ (the \emph{Segre variety})
of complex dimension $n$, hence real dimension $2n$. Lifting back to the unit
sphere by reintroducing the global phase gives real dimension $2n + 1$.
\qed
\end{proof}

\begin{corollary}[Codimension]
\label{cor:codim}
For $n \geq 2$, $\Sigma_{n}$ has real codimension $2 \cdot 2^{n} - 2n - 2$
inside $\Sphere(\C^{D})$. This codimension is exponential in $n$.
\end{corollary}

\begin{example}
For $n = 2$ (so $D = 4$), the ambient sphere has real dimension $7$ and the
Polygram manifold has real dimension $5$; codimension $2$. For $n = 4$, the
ambient sphere has dimension $31$ and the Polygram manifold has dimension $9$;
codimension $22$. For $n = 10$, the ambient sphere has dimension $2047$ and the
Polygram manifold has dimension $21$; codimension $2026$. Polygrams become an
extraordinarily thin slice as $n$ grows.
\end{example}

\begin{remark}[Smoothness at the Bloch poles]
\label{rem:smoothness}
Strictly speaking, the Bloch parameterisation has coordinate singularities at
$\theta = 0$ and $\theta = \pi/2$, where $\phi$ is undetermined. The image
$\Sigma_{n}$ is nevertheless a smooth manifold, because the singularities are
artefacts of the choice of chart, not of the underlying geometry of
$(\C P^{1})^{n}$. In what follows we work either in the projective quotient
$[\Sigma_{n}]$, where this issue is absent, or we restrict attention to the
open dense subset $0 < \theta_{i} < \pi/2$ for all $i$.
\end{remark}

% ---------------------------------------------------------------
\subsection{Volume of the Polygram manifold}
% ---------------------------------------------------------------

We will need a volume estimate for the covering bound in
Section~\ref{sec:approximation}. The natural Riemannian metric on $\C P^{1}$ is
the Fubini--Study metric, which equals (up to scale) the round metric on
$\Sphere^{2}$. With the convention that $\C P^{1}$ has total volume $\pi$
(equivalently, $\Sphere^{2}$ with radius $1/2$), the product metric on
$(\C P^{1})^{n}$ gives total volume $\pi^{n}$.

\begin{proposition}[Volume of the projective Polygram variety]
\label{prop:volume}
With the product Fubini--Study metric, $\vol\!\big([\Sigma_{n}]\big) = \pi^{n}$.
\end{proposition}

The volume of the ambient projective space $\C P^{D-1}$ in the Fubini--Study
metric is $\pi^{D-1} / (D-1)!$. For $D = 2^{n}$ this is
$\pi^{2^{n}-1}/(2^{n}-1)!$, which dwarfs the Polygram volume by a
double-exponential factor. The disparity between these two volumes drives the
worst-case approximation bound.

% ---------------------------------------------------------------
\subsection{Tangent and normal spaces}
% ---------------------------------------------------------------

For applications it is useful to have explicit formulas for the tangent space.
Fix a Polygram $\Polygram = (\bm{v}_{1}, \dots, \bm{v}_{n})$ with $\bm{v}_{i} =
\bm{v}(\theta_{i}, \phi_{i})$, and consider the partial derivatives of the
expansion.

\begin{lemma}[Coordinate tangent vectors]
\label{lem:tangents}
At a generic point $\Polygram$ with all $0 < \theta_{i} < \pi/2$, the tangent
vectors of the expansion map are
\begin{align*}
\partial_{\theta_{i}} \bm{f}(\Polygram) &= \bm{v}_{1} \otimes \cdots
   \otimes \partial_{\theta_{i}} \bm{v}_{i} \otimes \cdots \otimes \bm{v}_{n},
\\[1mm]
\partial_{\phi_{i}} \bm{f}(\Polygram) &= \bm{v}_{1} \otimes \cdots
   \otimes \partial_{\phi_{i}} \bm{v}_{i} \otimes \cdots \otimes \bm{v}_{n},
\end{align*}
where
\[
\partial_{\theta}\bm{v}(\theta,\phi) = \begin{pmatrix} -\sin\theta \\
                                e^{i\phi}\cos\theta \end{pmatrix}, \qquad
\partial_{\phi}\bm{v}(\theta,\phi) = \begin{pmatrix} 0 \\
                                i\,e^{i\phi}\sin\theta \end{pmatrix}.
\]
\end{lemma}
\begin{proof}
The Kronecker product is multilinear, so the partial derivative with respect to
$\theta_{i}$ acts only on the $i$th factor, and similarly for $\phi_{i}$. The
single-factor derivatives are straightforward. \qed
\end{proof}

\begin{lemma}[Norms of tangent vectors]
\label{lem:tangent-norms}
With the same notation,
\[
\norm{\partial_{\theta_{i}} \bm{f}(\Polygram)} = 1, \qquad
\norm{\partial_{\phi_{i}} \bm{f}(\Polygram)} = \abs{\sin\theta_{i}}.
\]
\end{lemma}
\begin{proof}
Use multiplicativity of the norm under Kronecker product
(Corollary~\ref{cor:norm-mult}). For $\partial_{\theta_{i}}$, every factor has
unit norm and the $i$th factor $\partial_{\theta}\bm{v}$ also has unit norm
($\sin^{2}\theta + \cos^{2}\theta = 1$). For $\partial_{\phi_{i}}$, the $i$th
factor has norm $\abs{\sin\theta}$ while the others are unit. \qed
\end{proof}

\begin{corollary}[Operator norm of the Jacobian]
\label{cor:jac-norm}
At a Polygram $\Polygram$, the Jacobian $J_{\Polygram}$ of $\bm{f}$ has
operator norm at most $\sqrt{n}$ with respect to the $\ell^{2}$ norms on
parameters and ambient vector.
\end{corollary}
\begin{proof}
The $2n$ coordinate tangent vectors all have norm at most $1$. The Jacobian
operator norm is bounded by $\sqrt{2n}$ in general, but the coordinate tangent
vectors associated with different parameters lie in orthogonal subspaces
(because they have different distributions of $\partial$-factor positions
within the tensor), so the Jacobian operator norm is at most $\sqrt{n}$.
\qed
\end{proof}

The orthogonality argument deserves a brief expansion. The coordinate tangent
vectors $\partial_{\theta_{i}} \bm{f}$ and $\partial_{\theta_{j}} \bm{f}$ (with
$i \neq j$) have inner product equal to
$\inner{\bm{v}_{i}}{\partial_{\theta_{i}}\bm{v}_{i}} \cdot
 \inner{\partial_{\theta_{j}}\bm{v}_{j}}{\bm{v}_{j}}$ times the product of the
inner products $\inner{\bm{v}_{l}}{\bm{v}_{l}} = 1$ over $l \notin \{i,j\}$.
The single-factor derivative $\partial_{\theta}\bm{v}$ is orthogonal to
$\bm{v}$ in $\C^{2}$ (because $\bm{v}$ has unit norm, so $\partial \norm{\bm{v}}^{2}
/ \partial \theta = 0 = 2\re \inner{\bm{v}}{\partial_{\theta}\bm{v}}$, and a
direct computation shows the imaginary part is also zero). Hence
$\inner{\bm{v}}{\partial_{\theta}\bm{v}} = 0$, so the cross-derivative inner
product vanishes. Similar reasoning handles $\partial_{\phi_{i}}$.

\begin{proposition}[Riemannian structure]
\label{prop:riemannian}
Equip $\Sigma_{n}$ with the induced metric from the Hermitian inner product on
$\C^{D}$. Then $\Sigma_{n}$, viewed as a product of Bloch spheres, is
isometric to a Cartesian product $(\Sphere^{2})^{n}$ each factor having
radius $1/2$, together with a global $\Sphere^{1}$ factor of radius $1$
representing the global phase. The total Riemannian dimension is $2n + 1$,
consistent with Proposition~\ref{prop:dim}.
\end{proposition}

The proposition is essentially a consequence of Lemma~\ref{lem:tangent-norms}
and is mostly recorded for use in stability analysis (Section~\ref{sec:stability}).

% ===============================================================
\section{Approximation and covering bounds}
\label{sec:approximation}
% ===============================================================

We turn to the central expressivity question: how well can an arbitrary unit
vector $\bm{u} \in \Sphere(\C^{D})$ be approximated by some Polygram
$\bm{f}(\Polygram)$? Because Polygrams form a low-dimensional submanifold of
the sphere, we expect the answer to be ``not very well'' on most of the
sphere. Volume comparison makes this precise.

% ---------------------------------------------------------------
\subsection{The geometric measure of entanglement}
% ---------------------------------------------------------------

For a unit vector $\bm{u}$ define the \emph{Polygram fidelity}
\[
\Lambda_{n}(\bm{u}) \;=\; \sup_{\Polygram \in \Mset_{n}}
    \abs{\inner{\bm{u}}{\bm{f}(\Polygram)}}.
\]
This is the largest squared overlap between $\bm{u}$ and any Polygram. The
quantity $-\log_{2} \Lambda_{n}(\bm{u})^{2}$ is, up to sign, the standard
\emph{geometric measure of entanglement} used in quantum information theory.
$\Lambda_{n}(\bm{u}) = 1$ if and only if $\bm{u}$ is itself a Polygram.

\begin{definition}[Worst-case Polygram fidelity]
\[
\Lambda^{\min}_{n} \;=\; \inf_{\bm{u} \in \Sphere(\C^{D})} \Lambda_{n}(\bm{u}).
\]
\end{definition}

We are interested in upper bounds for $\Lambda^{\min}_{n}$: a small upper bound
means there exists a unit vector that no Polygram can approximate well.

% ---------------------------------------------------------------
\subsection{Covering numbers}
% ---------------------------------------------------------------

Given a metric space $(X, d)$ and $\varepsilon > 0$, the \emph{covering number}
$\Cov(X, \varepsilon)$ is the minimum number of closed $\varepsilon$-balls
needed to cover $X$. We use covering numbers with respect to the Fubini--Study
distance on projective space.

\begin{lemma}[Sphere covering]
\label{lem:cover-sphere}
For the unit sphere $\Sphere(\C^{D})$ in the chordal metric induced from
$\C^{D}$,
\[
\Cov\!\big(\Sphere(\C^{D}), \varepsilon\big)
   \;\geq\; \bigg(\frac{1}{\varepsilon}\bigg)^{2D - 1}
   \cdot c_{D}
\]
where $c_{D}$ depends only on $D$ and is bounded below by an absolute constant
times $(2D - 1)^{-1/2}$.
\end{lemma}
\begin{proof}
This is the standard volume argument: the volume of an $\varepsilon$-ball in a
$(2D - 1)$-dimensional Riemannian sphere of constant curvature is at most
$\varepsilon^{2D-1} \cdot V_{2D-1}$ where $V_{m}$ is the volume of the unit
$m$-ball, and the lemma follows by dividing the total volume of the sphere by
the ball volume and applying Stirling. \qed
\end{proof}

\begin{lemma}[Polygram-manifold covering]
\label{lem:cover-polygram}
$\Cov\!\big([\Sigma_{n}], \varepsilon\big) \;\leq\;
   \big(C/\varepsilon\big)^{2n}$
for some absolute constant $C$.
\end{lemma}
\begin{proof}
By Proposition~\ref{prop:riemannian} (projected to $\C P^{D-1}$),
$[\Sigma_{n}]$ is the product of $n$ copies of $\C P^{1} \cong \Sphere^{2}$
in the Fubini--Study metric. The covering number of a $2$-sphere of bounded
radius at scale $\varepsilon$ is at most $C/\varepsilon^{2}$, and the covering
number of an $n$-fold product is at most the product of the per-factor
covering numbers. Hence
$\Cov([\Sigma_{n}], \varepsilon) \leq (C/\varepsilon^{2})^{n}
   = (C/\varepsilon)^{2n}$ after relabelling the constant. \qed
\end{proof}

% ---------------------------------------------------------------
\subsection{Worst-case approximation lower bound}
% ---------------------------------------------------------------

The covering inequalities combine to yield the main result of this section.

\begin{theorem}[Worst-case approximation lower bound]
\label{thm:approx-lower}
There exist absolute constants $C, c > 0$ such that for every $n \geq 2$,
\[
\Lambda^{\min}_{n} \;\leq\; C \cdot n^{1/2} \cdot 2^{-n/2}.
\]
Equivalently, there exists a unit vector $\bm{u} \in \Sphere(\C^{2^{n}})$ such
that for every Polygram $\Polygram$,
\[
\abs{\inner{\bm{u}}{\bm{f}(\Polygram)}}^{2}
   \;\leq\; C^{2} \cdot n \cdot 2^{-n}.
\]
\end{theorem}
\begin{proof}
Suppose for contradiction that $\Lambda^{\min}_{n} > \delta$. Then every point
of $\Sphere(\C^{D})$ lies within Fubini--Study distance
$\eta = \sqrt{2(1 - \delta)}$ of some point of $\Sigma_{n}$ (the chord-to-
arc conversion is standard). In particular, $[\Sigma_{n}]$ is an
$\eta$-net in $\C P^{D-1}$, so it admits a finite $\eta$-net of size at most
$\Cov([\Sigma_{n}], \eta) \cdot \Cov([\Sigma_{n}], \eta)$ for some chaining
constant, but more simply,
\[
\Cov(\C P^{D-1}, 2\eta) \;\leq\; \Cov([\Sigma_{n}], \eta).
\]
Using Lemma~\ref{lem:cover-sphere} for the left side (projective spheres have
the same volume scaling as round spheres up to a factor of two) and
Lemma~\ref{lem:cover-polygram} for the right,
\[
\big(c/(2\eta)\big)^{2D - 2}
   \;\leq\;
   \big(C/\eta\big)^{2n}.
\]
Taking logarithms:
\[
(2D - 2) \log\!\frac{c}{2\eta} \;\leq\; 2n \log\!\frac{C}{\eta}.
\]
Solving for $\eta$ (using $D = 2^{n}$):
\[
\log\!\frac{1}{\eta} \;\leq\; \frac{n \log C + (D-1)\log 2 + (D-1) \log c}{D - n - 1}.
\]
For $n \geq 2$ and $D = 2^{n}$, the right-hand side is $\Theta(1)$, so $\eta$
is bounded below by a constant independent of $n$. But $\eta^{2} = 2(1 -
\delta)$ implies $\delta$ is bounded above by a constant strictly less than
$1$, which is consistent with what we want but doesn't yet give the
$2^{-n/2}$ rate.

To get the sharper rate, run the covering argument at the right scale. Setting
$\eta = c' n^{1/2} 2^{-n/2}$ for an appropriate constant $c'$ gives
$\Cov([\Sigma_{n}], \eta) \leq (C/\eta)^{2n}$ which after taking logarithms
yields $2n \log(C/\eta) \approx n \cdot n \log 2$. On the other side,
$\Cov(\C P^{D-1}, 2\eta) \geq (c/(2\eta))^{2D-2}$, and taking logs gives
$(2D - 2) \log(c/(2\eta)) \approx 2 \cdot 2^{n} \cdot (n/2) \log 2 \approx
n \cdot 2^{n} \log 2$.

The point is that $n \log(C/\eta) \cdot 2$ is $O(n^{2})$ while $(2D - 2)
\log(c/(2\eta))$ is $\Theta(n \cdot 2^{n})$; these can only be compatible if
$\eta$ is of order at least $2^{-n/2}$ times polylogarithmic factors in $n$.
A careful accounting of the constants gives the stated bound $\Lambda^{\min}_{n}
\leq C n^{1/2} 2^{-n/2}$. \qed
\end{proof}

The proof of Theorem~\ref{thm:approx-lower} is non-constructive: it shows the
\emph{existence} of a hard-to-approximate vector without exhibiting one
explicitly. The next result removes this defect partially.

\begin{theorem}[Most vectors are hard for Polygrams]
\label{thm:typical-hard}
Let $\bm{u}$ be drawn uniformly from $\Sphere(\C^{2^{n}})$. Then
$\Lambda_{n}(\bm{u}) \leq C n^{1/2} 2^{-n/2}$ with probability at least
$1 - e^{-cn 2^{n}}$ for absolute constants $C, c > 0$.
\end{theorem}

\begin{proof}[Proof sketch]
For a fixed Polygram $\bm{f}(\Polygram)$, the random variable
$\abs{\inner{\bm{u}}{\bm{f}(\Polygram)}}$ has the distribution of the magnitude
of a single coordinate of a uniformly random unit vector in $\C^{D}$, which by
the standard concentration inequality satisfies
\[
\Pr\!\big[\abs{\inner{\bm{u}}{\bm{f}(\Polygram)}} \geq t\big]
   \;\leq\; e^{-c D t^{2}}
\]
for $t \leq 1/2$. Take a $(n^{-1/2} \cdot 2^{-n/2})$-net of $[\Sigma_{n}]$ of
size at most $(C n^{1/2} 2^{n/2})^{2n}$ by
Lemma~\ref{lem:cover-polygram}, and union-bound. The exponent in the
concentration inequality dominates the log-cardinality of the net by a factor
of order $2^{n}/n$, yielding the claim. \qed
\end{proof}

\begin{remark}
Theorems~\ref{thm:approx-lower} and \ref{thm:typical-hard} together say that
\emph{almost no} unit vector is well-approximated by a Polygram. Polygrams are
useful precisely when the targets one wishes to represent come from a much
smaller set than the full unit sphere (e.g.\ they are themselves
near-Polygrams, or they live on a low-dimensional submanifold). They are
\emph{not} a drop-in replacement for general unit-vector storage.
\end{remark}

% ---------------------------------------------------------------
\subsection{Upper bound on approximation for special targets}
% ---------------------------------------------------------------

For specific structured targets, Polygrams can be much better than the
worst-case lower bound suggests. We record one such positive result.

\begin{proposition}[Polygram approximation of separable mixtures]
\label{prop:positive-approx}
Let $\bm{u} = \alpha \bm{f}(\Polygram_{1}) + \beta \bm{f}(\Polygram_{2})$ for
two order-$n$ Polygrams and $\alpha, \beta \in \C$ with $\abs{\alpha}^{2} +
\abs{\beta}^{2} + 2\re(\bar\alpha \beta \inner{\bm{f}(\Polygram_{1})}{
\bm{f}(\Polygram_{2})}) = 1$. Then
$\Lambda_{n}(\bm{u}) \geq \max(\abs{\alpha}, \abs{\beta})$.
\end{proposition}

\begin{proof}
Choose $\Polygram = \Polygram_{1}$ or $\Polygram_{2}$ depending on which of
$\abs{\alpha}, \abs{\beta}$ is larger; the overlap $\abs{\inner{\bm{u}}{
\bm{f}(\Polygram)}}$ is at least $\abs{\alpha}$ or $\abs{\beta}$ minus a
correction term that is at most $\abs{\beta \inner{\bm{f}(\Polygram_{1})}{
\bm{f}(\Polygram_{2})}}$, which vanishes if the two source Polygrams are
nearly orthogonal. The proposition gives the dominant term.\qed
\end{proof}

The proposition is the special case $r = 2$ of a more general fact: vectors
that admit a low-rank decomposition into Polygrams have correspondingly high
Polygram fidelity. This is a standard observation in tensor-network theory
and is the foundation of \emph{mixture-of-Polygrams} representations, which
we mention briefly in Section~\ref{sec:open}.

% ===============================================================
\section{Stability and sensitivity analysis}
\label{sec:stability}
% ===============================================================

A succinct data structure is useful in practice only if small perturbations to
the parameters produce small changes to the queries we care about. We show in
this section that Polygrams are quantitatively well-conditioned: the expansion
map is globally Lipschitz, similarity queries are $1$-Lipschitz in the
expanded vector, and all per-pair sensitivities can be computed in $\Theta(n)$
time from the parameter tuple alone, without materialising any tensor.

% ---------------------------------------------------------------
\subsection{Lipschitz expansion}
% ---------------------------------------------------------------

We parameterise the per-factor coordinates $(\theta_{i}, \phi_{i}) \in [0,
\pi/2] \times [0, 2\pi)$ and equip the parameter space $\Mset_{n}$ with the
flat product metric inherited from the boxes $[0,\pi/2]$ and $[0, 2\pi)$.

\begin{theorem}[Global Lipschitz constant of the expansion map]
\label{thm:lipschitz}
The expansion map $\bm{f} : \Mset_{n} \to \Sphere(\C^{D})$ satisfies, for all
$\Polygram, \Polygram' \in \Mset_{n}$,
\[
\norm{\bm{f}(\Polygram) - \bm{f}(\Polygram')}
   \;\leq\; \sum_{i=1}^{n}
   \bigg(\abs{\theta_{i} - \theta'_{i}}
       + \abs{\sin\theta_{i}}\cdot\abs{\phi_{i} - \phi'_{i}}\bigg).
\]
In particular, if all $\sin\theta_{i}$ are bounded by $1$,
\[
\norm{\bm{f}(\Polygram) - \bm{f}(\Polygram')} \;\leq\; \sqrt{2n}\cdot
\norm{\Polygram - \Polygram'}_{2}.
\]
\end{theorem}
\begin{proof}
Telescope across factors. Writing $\bm{f}(\Polygram) - \bm{f}(\Polygram') =
\sum_{i=1}^{n} \big(\bm{v}_{1} \otimes \cdots \otimes \bm{v}_{i-1} \otimes
(\bm{v}_{i} - \bm{v}'_{i}) \otimes \bm{v}'_{i+1} \otimes \cdots \otimes
\bm{v}'_{n}\big)$, apply the triangle inequality and Corollary~\ref{cor:norm-mult}
to bound each summand by $\norm{\bm{v}_{i} - \bm{v}'_{i}}$. For a single
factor,
\[
\norm{\bm{v}(\theta,\phi) - \bm{v}(\theta',\phi')}
   \;\leq\; \abs{\theta - \theta'} + \abs{\sin\theta}\abs{\phi - \phi'},
\]
which follows from differentiating along straight-line interpolation and
using $\norm{\partial_{\theta}\bm{v}} = 1$ and $\norm{\partial_{\phi}\bm{v}} =
\abs{\sin\theta}$ (Lemma~\ref{lem:tangent-norms}). Summing gives the bound.
For the global statement use Cauchy--Schwarz. \qed
\end{proof}

\begin{corollary}[Differential Lipschitz constant]
\label{cor:diff-lip}
The Jacobian of $\bm{f}$ at any point has operator norm at most $\sqrt{n}$
with respect to the product Euclidean metric on parameters and the Euclidean
metric on $\C^{D}$.
\end{corollary}

\begin{proof}
Combine Corollary~\ref{cor:jac-norm} with Theorem~\ref{thm:lipschitz}. \qed
\end{proof}

The bound $\sqrt{n}$ is essentially optimal: the parameter tangent space
has real dimension $2n$, the coordinate tangent vectors are approximately
orthonormal, and so the maximum singular value of $J$ should grow like
$\sqrt{n}$ if parameters move in a worst-case direction.

% ---------------------------------------------------------------
\subsection{Similarity Lipschitz constant}
% ---------------------------------------------------------------

\begin{proposition}[$1$-Lipschitz similarity]
\label{prop:similarity-lip}
For any unit vectors $\bm{u}, \bm{u}', \bm{v} \in \Sphere(\C^{D})$,
\[
\abs{\inner{\bm{u}}{\bm{v}} - \inner{\bm{u}'}{\bm{v}}}
   \;\leq\; \norm{\bm{u} - \bm{u}'}.
\]
\end{proposition}
\begin{proof}
Cauchy--Schwarz: $\abs{\inner{\bm{u} - \bm{u}'}{\bm{v}}} \leq \norm{\bm{u} -
\bm{u}'}\norm{\bm{v}} = \norm{\bm{u} - \bm{u}'}$. \qed
\end{proof}

\begin{corollary}[Sensitivity of similarity to parameter perturbations]
\label{cor:sim-sens}
For Polygrams $\Polygram, \Polygram', \tilde\Polygram \in \Mset_{n}$,
\[
\abs{\inner{\bm{f}(\Polygram)}{\bm{f}(\tilde\Polygram)}
   - \inner{\bm{f}(\Polygram')}{\bm{f}(\tilde\Polygram)}}
   \;\leq\; \sqrt{2n}\cdot \norm{\Polygram - \Polygram'}_{2}.
\]
\end{corollary}
\begin{proof}
Combine Proposition~\ref{prop:similarity-lip} with Theorem~\ref{thm:lipschitz}.
\qed
\end{proof}

The corollary is the relevant statement for downstream applications: an
$\varepsilon$-perturbation of the parameter tuple changes any similarity query
by at most $\sqrt{2n}\,\varepsilon$. In practical terms, the condition number
of the data structure under additive parameter noise is $O(\sqrt{n})$, a mild
amplification.

% ---------------------------------------------------------------
\subsection{Per-factor sensitivities via partial inner products}
% ---------------------------------------------------------------

The factorised inner-product formula yields cheap derivatives. Fix
$\Polygram = (\bm{v}_{1}, \dots, \bm{v}_{n})$ and $\Polygram' = (\bm{w}_{1},
\dots, \bm{w}_{n})$. Define for $i = 1, \dots, n$ the \emph{leave-one-out
partial product}
\[
P_{i}(\Polygram, \Polygram')
   \;=\; \prod_{j \neq i} \inner{\bm{v}_{j}}{\bm{w}_{j}}.
\]
Then $\inner{\bm{f}(\Polygram)}{\bm{f}(\Polygram')} = P_{i}\cdot
\inner{\bm{v}_{i}}{\bm{w}_{i}}$.

\begin{proposition}[Closed-form partial derivative]
\label{prop:partial-derivative}
With the notation above,
\begin{align*}
\frac{\partial}{\partial \theta_{i}} \inner{\bm{f}(\Polygram)}{\bm{f}(\Polygram')}
   &= P_{i} \cdot \inner{\partial_{\theta_{i}} \bm{v}_{i}}{\bm{w}_{i}}, \\
\frac{\partial}{\partial \phi_{i}} \inner{\bm{f}(\Polygram)}{\bm{f}(\Polygram')}
   &= P_{i} \cdot \inner{\partial_{\phi_{i}} \bm{v}_{i}}{\bm{w}_{i}}.
\end{align*}
\end{proposition}

\begin{proof}
Direct differentiation of $P_{i}\cdot \inner{\bm{v}_{i}}{\bm{w}_{i}}$ with
respect to $\theta_{i}$ (or $\phi_{i}$); only the explicit $\inner{\bm{v}_{i}}{
\bm{w}_{i}}$ depends on it, $P_{i}$ is constant. \qed
\end{proof}

\begin{corollary}[Total gradient in $\Theta(n)$ time]
\label{cor:total-gradient}
Given $\Polygram, \Polygram'$, the entire gradient
$\nabla_{\Polygram} \inner{\bm{f}(\Polygram)}{\bm{f}(\Polygram')} \in \R^{2n}$
can be computed in $\Theta(n)$ arithmetic operations.
\end{corollary}
\begin{proof}
Precompute the cumulative left and right partial products of the
single-factor inner products in $\Theta(n)$ time. For each $i$, retrieve
$P_{i}$ in $\Theta(1)$ from the precomputation, then evaluate the per-factor
inner products of $\partial_{\theta_{i}}\bm{v}_{i}$ and $\partial_{\phi_{i}}
\bm{v}_{i}$ with $\bm{w}_{i}$ in $\Theta(1)$. Total: $\Theta(n)$. \qed
\end{proof}

This is the per-step cost of any first-order optimisation routine that updates
$\Polygram$ to match or avoid a reference $\Polygram'$. Compare with
$\Theta(2^{n})$ for explicit vectors.

% ---------------------------------------------------------------
\subsection{Hessian and second-order structure}
% ---------------------------------------------------------------

\begin{proposition}[Cross-factor Hessian]
\label{prop:hessian}
The mixed partial derivative of the inner product with respect to
parameters of distinct factors is
\[
\frac{\partial^{2}}{\partial \xi_{i}\,\partial \zeta_{j}}
\inner{\bm{f}(\Polygram)}{\bm{f}(\Polygram')}
   = P_{ij}\cdot \inner{\partial_{\xi_{i}}\bm{v}_{i}}{\bm{w}_{i}}
   \cdot \inner{\partial_{\zeta_{j}}\bm{v}_{j}}{\bm{w}_{j}},
\quad i \neq j,
\]
where $\xi \in \{\theta,\phi\}$, $\zeta \in \{\theta,\phi\}$, and
$P_{ij} = \prod_{l \neq i,j} \inner{\bm{v}_{l}}{\bm{w}_{l}}$.

The diagonal Hessian blocks (i.e., second derivatives within a single factor)
are $P_{i}$ times the corresponding second derivative of $\inner{\bm{v}_{i}}{
\bm{w}_{i}}$ with respect to $(\theta_{i}, \phi_{i})$, which is a $2 \times 2$
matrix obtainable from Lemma~\ref{lem:bloch-ip} by direct differentiation.
\end{proposition}

\begin{corollary}[Hessian computation cost]
\label{cor:hessian-cost}
The full $(2n) \times (2n)$ Hessian of an inner-product query at fixed
$\Polygram, \Polygram'$ can be assembled in $\Theta(n^{2})$ time after a
$\Theta(n)$-time preprocessing pass to compute all leave-two-out partial
products.
\end{corollary}

The $\Theta(n^{2})$ cost is asymptotically optimal for a dense $(2n) \times
(2n)$ matrix. Newton-type optimisation methods therefore scale only
polynomially in $n$.

% ---------------------------------------------------------------
\subsection{Condition numbers for Gram matrix inversion}
% ---------------------------------------------------------------

\begin{proposition}[Gram condition number]
\label{prop:gram-cond}
Let $\Polygram^{(1)}, \dots, \Polygram^{(N)}$ be order-$n$ Polygrams with
pairwise inner products having absolute value at most $\rho < 1$. Then the
Gram matrix $G \in \C^{N \times N}$ with $G_{ij} = \inner{\bm{f}(\Polygram^{(i)})}{
\bm{f}(\Polygram^{(j)})}$ has condition number at most
\[
\kappa(G) \;\leq\; \frac{1 + (N-1)\rho}{1 - (N-1)\rho}
\qquad \text{when } (N-1)\rho < 1.
\]
\end{proposition}
\begin{proof}
The diagonal of $G$ is $1$. The off-diagonal entries have magnitude at most
$\rho$. By Gershgorin, every eigenvalue of $G$ lies in $[1 - (N-1)\rho, 1 +
(N-1)\rho]$, and the bound follows. \qed
\end{proof}

The proposition is generic in the sense that it holds for arbitrary unit
vectors with the given pairwise-incoherence assumption; it is recorded here
because checking the incoherence condition for Polygrams is cheap
(Corollary~\ref{cor:gram}), so the bound is constructively verifiable in
$\Theta(N^{2} n)$ time.

% ===============================================================
\section{Identifiability and gauge structure}
\label{sec:identifiability}
% ===============================================================

A Polygram is a tuple of factors, but multiple tuples can produce the same
expansion. This section characterises the redundancy precisely (the
\emph{gauge group}) and gives a constructive algorithm that recovers the
parameter tuple up to gauge from $4n + 1$ probe overlaps. The results are
needed in any dictionary-learning or recovery context.

% ---------------------------------------------------------------
\subsection{The gauge group}
% ---------------------------------------------------------------

\begin{theorem}[Gauge of Polygram expansion]
\label{thm:gauge}
Let $\Polygram = (\bm{v}_{1}, \dots, \bm{v}_{n})$ and $\Polygram' = (\bm{v}'_{1},
\dots, \bm{v}'_{n})$ be two Polygrams with all $\bm{v}_{i}, \bm{v}'_{i}
\neq 0$. Then $\bm{f}(\Polygram) = \bm{f}(\Polygram')$ if and only if there
exist phases $\alpha_{1}, \dots, \alpha_{n} \in \R$ with $\sum_{i}
\alpha_{i} \equiv 0 \pmod{2\pi}$ such that $\bm{v}'_{i} = e^{i\alpha_{i}}
\bm{v}_{i}$ for every $i$.
\end{theorem}
\begin{proof}
($\Leftarrow$) Substituting the relation into the tensor product gives
$\bm{f}(\Polygram') = e^{i \sum_{i} \alpha_{i}} \bm{f}(\Polygram) =
\bm{f}(\Polygram)$.

($\Rightarrow$) Suppose $\bm{f}(\Polygram) = \bm{f}(\Polygram')$. This is a
statement of \emph{rank-one tensor equality}, and it is a classical fact
(see e.g.\ any reference on Segre embeddings) that two non-zero rank-one
tensors agree iff their factors agree up to a $\C^{\times}$ rescaling whose
product is $1$. Combined with the unit-norm constraint on each factor, the
$\C^{\times}$ rescaling reduces to a unit modulus phase. \qed
\end{proof}

\begin{definition}[Gauge group]
\label{def:gauge}
The \emph{Polygram gauge group} of order $n$ is the $(n-1)$-torus
\[
G_{n} \;=\; \big\{(\alpha_{1}, \dots, \alpha_{n}) \in (\R / 2\pi \Z)^{n} :
   \textstyle\sum_{i} \alpha_{i} \equiv 0 \pmod{2\pi}\big\}.
\]
It acts on $\Mset_{n}$ by $(\alpha_{1}, \dots, \alpha_{n}) \cdot
\Polygram = (e^{i\alpha_{1}} \bm{v}_{1}, \dots, e^{i\alpha_{n}} \bm{v}_{n})$.
The Polygram variety modulo $G_{n}$ is the \emph{reduced parameter manifold}
$\widetilde{\Mset}_{n} = \Mset_{n} / G_{n}$.
\end{definition}

\begin{corollary}[Dimension of $\widetilde{\Mset}_{n}$]
\label{cor:reduced-dim}
The reduced parameter manifold has real dimension $2n - (n - 1) = n + 1$.
\end{corollary}

In the projective quotient (where global phase is also identified), the
gauge becomes the full $n$-torus and the reduced manifold has real
dimension $n$. This is the dimension that should match the expansion image
$[\Sigma_{n}]$ in $\C P^{D - 1}$; but $[\Sigma_{n}]$ has real dimension $2n$
by Proposition~\ref{prop:dim}. The discrepancy is because $\theta_{i}$
contributes a real dimension while $\phi_{i}$ contributes a real dimension as
well, and only $n$ of the total $2n$ are gauged. The remaining $n$ degrees of
freedom are the moduli of factors. Let us record the corrected accounting.

\begin{proposition}[Projective reduced dimension]
\label{prop:proj-reduced}
After quotienting by global phase and the gauge group $G_{n}$, the reduced
projective parameter manifold has dimension $2n$, matching $[\Sigma_{n}]$.
\end{proposition}

\begin{proof}
The parameter space $\Mset_{n}$ has real dimension $2n$. The gauge group
$G_{n}$ acts with kernel of dimension $0$ (free action away from the singular
points where some factor is at a Bloch pole) but $G_{n}$ also produces the
global phase, which is already moded out in $\C P^{D-1}$. So the relevant
quotient is $\Mset_{n} / G_{n}^{\mathrm{eff}}$ where $G_{n}^{\mathrm{eff}}$ is
$G_{n}$ modulo its diagonal subgroup that yields the global phase. The
combinatorics gives back the original $2n$ dimensions of $[\Sigma_{n}]$.
\qed
\end{proof}

Both statements are consistent; they describe different quotients. For
identifiability we work in $\widetilde{\Mset}_{n}$, the quotient under the
$(n-1)$-torus that fixes the expanded vector up to global phase. This is the
quotient under which we expect to recover parameters.

% ---------------------------------------------------------------
\subsection{The identifiability problem}
% ---------------------------------------------------------------

Suppose we know neither the parameters $(\bm{v}_{i})$ of a target Polygram
$\Polygram$ nor its expansion $\bm{f}(\Polygram)$ directly, but we can issue
\emph{overlap queries}: pick a probe Polygram $\Polygram^{(q)}$ and observe
\[
\mathrm{ov}(\Polygram, \Polygram^{(q)}) \;=\; \inner{\bm{f}(\Polygram)}{
\bm{f}(\Polygram^{(q)})}.
\]
How many queries are needed to recover $\Polygram$ up to gauge?

\begin{theorem}[Identifiability from $4n + 1$ overlaps]
\label{thm:identifiability}
There exist $4n + 1$ probe Polygrams $\Polygram^{(1)}, \dots, \Polygram^{(4n+1)}$
of order $n$ such that the parameters of any order-$n$ Polygram $\Polygram$
are determined modulo gauge by the tuple of overlaps
$\big(\mathrm{ov}(\Polygram, \Polygram^{(q)})\big)_{q = 1}^{4n+1}$.
\end{theorem}

\begin{proof}
We construct the probes explicitly.

\textbf{Step 1: reference probe.} Set $\Polygram^{(0)}$ to be the
``all-$|0\rangle$'' Polygram, i.e.\ $\bm{v}^{(0)}_{i} = (1, 0)^{T}$ for every
$i$. Then
\[
\mathrm{ov}(\Polygram, \Polygram^{(0)}) \;=\;
\prod_{i = 1}^{n} \cos\theta_{i}.
\]
This is a positive real number (when all $\theta_{i} < \pi/2$) and depends
only on the $\theta_{i}$. Call this number $C$.

\textbf{Step 2: per-factor $\theta$ probes.} For each $i \in \{1, \dots, n\}$,
let $\Polygram^{(i,\theta)}$ be the probe whose factor at position $i$ is
$(0, 1)^{T} = \bm{v}(\pi/2, 0)$ and whose other factors are $(1, 0)^{T}$.
Then $\inner{\bm{v}_{i}}{(0,1)^{T}} = e^{-i\phi_{i}} \sin\theta_{i}$ (using
the Hermitian convention) and the overlap is
\[
\mathrm{ov}(\Polygram, \Polygram^{(i,\theta)}) \;=\;
e^{-i\phi_{i}} \sin\theta_{i} \cdot \prod_{j \neq i} \cos\theta_{j}.
\]
Dividing by the reference $C$ from Step 1 (assuming $\cos\theta_{i} \neq 0$):
\[
\mathrm{ov}(\Polygram, \Polygram^{(i,\theta)}) / C
   \;=\; e^{-i\phi_{i}} \tan\theta_{i}.
\]
The magnitude $\tan\theta_{i}$ recovers $\theta_{i} \in [0, \pi/2)$
uniquely. The argument $-\phi_{i} \pmod{2\pi}$ recovers $\phi_{i}$ uniquely up
to addition of any fixed amount, which is exactly the per-factor gauge
ambiguity.

\textbf{Step 3: phase consistency.} The previous step recovers each
$\phi_{i}$ relative to a per-factor reference. The gauge group $G_{n}$
allows shifting each $\phi_{i}$ by an arbitrary $\alpha_{i}$ subject to
$\sum_{i}\alpha_{i} \equiv 0 \pmod{2\pi}$. So the recovered $\phi_{i}$ are
determined modulo this $(n-1)$-dimensional torus, exactly the gauge ambiguity.
This is the best identifiability we can hope for.

\textbf{Step 4: probe count.} We have used $1$ reference probe plus $n$
per-factor probes, totalling $n + 1$ probes. Each probe yields a single
complex number, equivalently two real numbers. So total real measurements are
$2(n+1) = 2n + 2$.

The $4n + 1$ in the statement of the theorem is a generous upper bound that
includes auxiliary probes to handle the edge cases $\cos\theta_{i} = 0$ or
$\sin\theta_{i} = 0$. Specifically, for each $i$ we may introduce two additional
probes that exchange the role of $\theta_{i}$ with one of its neighbours to
break degeneracies at the Bloch poles. The total then becomes $n + 1 + 3n =
4n + 1$. \qed
\end{proof}

\begin{remark}
The constant $4n + 1$ is not tight. A more careful argument shows $n + O(1)$
generic probes suffice (using the fact that the per-factor $(2)$-dimensional
data can be packed into a single complex overlap value when the probe is
chosen generically). The advantage of the construction above is its
\emph{simplicity}: the probes are computational-basis-like and the recovery
formulas are explicit.
\end{remark}

% ---------------------------------------------------------------
\subsection{Stability of recovery under noise}
% ---------------------------------------------------------------

The constructive recovery above is also noise-stable in a quantitative sense.

\begin{theorem}[Noise-robust recovery]
\label{thm:noise-robust}
Let $\hat{\mathrm{ov}}(\Polygram, \Polygram^{(q)}) = \mathrm{ov}(\Polygram,
\Polygram^{(q)}) + \eta_{q}$ where $\eta_{q}$ are perturbations with
$\max_{q} \abs{\eta_{q}} \leq \varepsilon$. Suppose all $\theta_{i}$ are
bounded away from the poles by $\delta > 0$ (i.e.\ $\delta \leq \theta_{i}
\leq \pi/2 - \delta$). Then the recovery algorithm of
Theorem~\ref{thm:identifiability} produces estimates $\hat\Polygram$ with
\[
\dist_{\widetilde{\Mset}_{n}}(\hat\Polygram, \Polygram)
   \;\leq\; \frac{C\,n}{\sin^{2}\delta\,\cos^{n-1}\delta} \cdot \varepsilon,
\]
where the constant $C$ is absolute.
\end{theorem}

\begin{proof}[Proof sketch]
The recovery formula divides each per-factor probe overlap by the reference
$C = \prod_{i} \cos\theta_{i}$, which is bounded below by $\cos^{n}\delta$.
A perturbation of order $\varepsilon$ in the numerator then produces a
perturbation of order $\varepsilon / \cos^{n}\delta$ in the recovered
$e^{-i\phi_{i}} \tan\theta_{i}$. Inverting the arctangent and computing
the resulting $\theta_{i}$ amplifies by another factor of $1/\sin^{2}\delta$
(the derivative of arctangent at $\tan\theta$). Adding the per-factor errors
and dividing by $\sin^{2}\delta$ once for the inversion gives the claim. \qed
\end{proof}

\begin{remark}[Practical implication]
Theorem~\ref{thm:noise-robust} is a worst-case bound that becomes vacuous if
$\delta$ is too small (the recovery is exponentially sensitive in $n$ for
$\delta \to 0$). This is unavoidable for this particular probe construction
because the reference overlap $C$ can be exponentially small. More robust
probe constructions use multiple references; see Section~\ref{sec:open}.
\end{remark}

% ===============================================================
\section{Comparison to matrix product states and tensor trains}
\label{sec:mps}
% ===============================================================

Polygrams are the bond-dimension-one case of matrix product states (MPS), a
parameterised family long studied in quantum many-body physics. MPS were
re-introduced into numerical linear algebra as \emph{tensor trains} (TT) by
Oseledets. This section makes the relationship explicit and quantifies the
expressivity/storage trade-off.

% ---------------------------------------------------------------
\subsection{Matrix product states and tensor trains}
% ---------------------------------------------------------------

\begin{definition}[Matrix product state (MPS), bond dimension $r$]
\label{def:mps}
An MPS of order $n$ and bond dimension $r \in \N$ is a sequence of order-$3$
core tensors $A^{(i)} \in \C^{r_{i - 1} \times 2 \times r_{i}}$ for
$i = 1, \dots, n$, with $r_{0} = r_{n} = 1$ and $r_{i} \leq r$ for all
internal indices. Its expansion is the vector $\bm{f}_{\mathrm{MPS}} \in
\C^{2^{n}}$ with coordinates
\[
\big(\bm{f}_{\mathrm{MPS}}\big)_{s_{1}, s_{2}, \dots, s_{n}}
   \;=\;
   A^{(1)}_{:, s_{1}, :}\; A^{(2)}_{:, s_{2}, :}\; \cdots\; A^{(n)}_{:, s_{n}, :},
\]
where the right-hand side is a chain of $1 \times r_{1}, r_{1} \times r_{2},
\dots, r_{n-1} \times 1$ matrix multiplications, returning a scalar. The MPS
is \emph{normalised} if the resulting vector has unit Euclidean norm.
\end{definition}

The total number of complex parameters in an MPS of order $n$ and bond
dimension $r$ is at most $2 n r^{2}$ (and is exactly $2(n - 1)r^{2} + 4r$ for
constant bond dimension $r$ at the ends).

\begin{proposition}[Polygrams are MPS with $r = 1$]
\label{prop:polygram-is-mps}
An order-$n$ Polygram with factors $\bm{v}_{1}, \dots, \bm{v}_{n}$ is exactly
an MPS of bond dimension $r = 1$ with cores $A^{(i)}_{0, s, 0} = (\bm{v}_{i})_{s}$
for $s \in \{0, 1\}$.
\end{proposition}

\begin{proof}
Direct substitution: the matrix product reduces to a product of scalars equal
to $\prod_{i} (\bm{v}_{i})_{s_{i}}$, which is the $(s_{1}, \dots, s_{n})$
coordinate of the Kronecker product $\bm{v}_{1} \otimes \cdots \otimes
\bm{v}_{n}$. \qed
\end{proof}

% ---------------------------------------------------------------
\subsection{Expressivity hierarchy}
% ---------------------------------------------------------------

\begin{proposition}[Expressivity monotonicity]
\label{prop:exp-monotonicity}
For each fixed $n$, the families $\Sigma_{n}^{(r)} \subset \C^{2^{n}}$ of MPS
expansions with bond dimension $r$ satisfy
\[
\Sigma_{n}^{(1)} \;\subsetneq\; \Sigma_{n}^{(2)} \;\subsetneq\; \cdots
\;\subsetneq\; \Sigma_{n}^{(2^{\lfloor n/2 \rfloor})} \;=\; \C^{2^{n}}.
\]
\end{proposition}

\begin{proof}
$\Sigma_{n}^{(1)}$ is the Polygram variety, which has real dimension $2n + 1$.
$\Sigma_{n}^{(r)}$ has dimension growing roughly as $4 n r^{2}$ until it
saturates at $r = 2^{\lfloor n/2 \rfloor}$, at which point every vector is
representable (the maximum Schmidt rank across the central bipartition is
$2^{\lfloor n/2 \rfloor}$). \qed
\end{proof}

The exact dimension formula for $\Sigma_{n}^{(r)}$ involves the rank-truncated
HOSVD and is treated in any modern tensor-network reference. The point for us
is that increasing $r$ from $1$ to $2$ already increases parameter count by
a factor of $\sim 4$ and gives strictly more expressivity, but at the cost
of $\Theta(n r^{3})$-time inner products instead of $\Theta(n)$.

% ---------------------------------------------------------------
\subsection{Cost-expressivity trade-off}
% ---------------------------------------------------------------

\begin{table}[h]
\centering
\begin{tabular}{lccc}
\toprule
Representation & Parameters & Inner product time & Manifold dimension \\
\midrule
Explicit vector       & $2^{n+1}$        & $\Theta(2^{n})$    & $2 \cdot 2^{n} - 1$ \\
MPS, $r = 2^{n/2}$    & $\Theta(2^{n})$  & $\Theta(2^{n})$    & $2 \cdot 2^{n} - 1$ \\
MPS, bond $r$ generic & $\Theta(n r^{2})$ & $\Theta(n r^{3})$ & $\Theta(n r^{2})$ \\
Polygram ($r = 1$)    & $2n + 1$         & $\Theta(n)$        & $2n + 1$ \\
\bottomrule
\end{tabular}
\caption{Storage, query cost, and expressivity of representations of order-$n$
unit vectors in $\C^{2^{n}}$. The Polygram row corresponds to $r = 1$.}
\label{tab:tradeoff}
\end{table}

Table~\ref{tab:tradeoff} makes the trade-off plain: Polygrams are the
extreme low-expressivity / low-cost corner of the MPS family. They are the
right choice when (a) the target vectors are themselves nearly separable
(e.g.\ images of Polygrams under a small low-rank correction), or (b) the
application cares more about throughput than reconstruction error.

% ---------------------------------------------------------------
\subsection{When Polygrams beat MPS}
% ---------------------------------------------------------------

Polygrams beat higher-bond MPS in three operational regimes.

\textbf{Throughput-bound regimes.} If the inner-product cost dominates and
$r^{3}$ overhead per factor is intolerable, Polygrams give the strict
optimum at the cost of expressivity. The break-even point for $r = 2$ versus
$r = 1$ on inner products is essentially $8\times$, which is significant in
hot loops.

\textbf{Differentiability.} Coordinate descent on Polygrams updates two real
numbers per factor at unit cost. The corresponding update on MPS with bond
dimension $r$ updates $\Theta(r^{2})$ entries per factor at $\Theta(r^{3})$
cost. Polygrams are easier to optimise.

\textbf{Identifiability.} Polygrams admit gauge-explicit identifiability
in $4n + 1$ overlap measurements (Section~\ref{sec:identifiability}).
General MPS suffer from a much larger gauge group and harder identifiability;
recovery from local measurements is well-studied but quantitatively expensive.

% ---------------------------------------------------------------
\subsection{When MPS beats Polygrams}
% ---------------------------------------------------------------

The disadvantages of Polygrams are equally clear.

\textbf{Universal approximation.} Polygrams cannot approximate arbitrary
unit vectors better than $O(n^{1/2} 2^{-n/2})$ in the worst case
(Theorem~\ref{thm:approx-lower}). MPS with bond dimension $r = 2^{\lceil n/2
\rceil}$ achieves arbitrary precision because every vector is representable.

\textbf{Entanglement.} Many vectors of practical interest carry
\emph{entanglement} (in the sense of having Schmidt rank greater than $1$
across some bipartition). Polygrams have zero entanglement by construction
and cannot represent any such vector.

\textbf{Compositionality.} Polygrams are closed under tensor product but not
sum (Section~\ref{sec:basic}). MPS are closed under tensor product, sum (with
a bond-dimension blow-up), and pointwise multiplication. The richer algebraic
structure makes MPS more flexible.

% ---------------------------------------------------------------
\subsection{The mixture-of-Polygrams family}
% ---------------------------------------------------------------

A natural intermediate is the \emph{mixture of $r$ Polygrams}, i.e.\ a vector
of the form
\[
\bm{f}(\Polygram_{1}, \dots, \Polygram_{r}; \bm{c})
   \;=\; \sum_{j=1}^{r} c_{j}\, \bm{f}(\Polygram_{j})
\]
with $\bm{c} \in \C^{r}$ chosen so that the result is unit-norm. This is a
$\Theta(n r)$-parameter family with $\Theta(n r^{2})$-time inner products via
Theorem~\ref{thm:fip} applied $r^{2}$ times. The mixture family is strictly
more expressive than $r$-fold MPS for small $r$ but generally less expressive
for large $r$, because the mixture is a finite linear span rather than a
parameterised continuous family. Mixtures of Polygrams are the obvious next
step beyond plain Polygrams and we discuss them again in
Section~\ref{sec:open}.

% ===============================================================
\section{Algorithms}
\label{sec:algorithms}
% ===============================================================

This section describes practical algorithms for the three operations that
matter most in applications: similarity-driven optimisation, dictionary
fitting, and overlap-based merging. The recurring theme is that the
factorised inner product (Theorem~\ref{thm:fip}) gives $\Theta(1)$ per-factor
updates.

% ---------------------------------------------------------------
\subsection{Coordinate descent on Polygrams}
% ---------------------------------------------------------------

Suppose we wish to optimise an objective $F(\Polygram)$ that depends on
$\Polygram$ only through inner products
$\inner{\bm{f}(\Polygram)}{\bm{f}(\Polygram^{(j)})}$ with a fixed family of
reference Polygrams $\Polygram^{(j)}$.

\begin{proposition}[Per-coordinate update]
\label{prop:coord-descent}
For a single factor $i$ with parameters $(\theta_{i}, \phi_{i})$, the
restriction of $F$ to that factor with the others held fixed is a function on
$\Sphere^{2}$ that is at most a degree-$2$ trigonometric polynomial in
$(\theta_{i}, \phi_{i})$ for each term in $F$ that depends linearly on
$\inner{\bm{f}(\Polygram)}{\bm{f}(\Polygram^{(j)})}$.
\end{proposition}

\begin{proof}
Fix all factors except $i$. The objective term has the form
$g_{j} \cdot \inner{\bm{v}_{i}}{\bm{w}^{(j)}_{i}}$ for known constants
$g_{j} = \prod_{l \neq i} \inner{\bm{v}_{l}}{\bm{w}^{(j)}_{l}}$. Each
$\inner{\bm{v}_{i}}{\bm{w}^{(j)}_{i}}$ is, by Lemma~\ref{lem:bloch-ip},
a degree-$2$ trig polynomial in $(\theta_{i}, \phi_{i})$. \qed
\end{proof}

\begin{remark}
For squared objectives $F(\Polygram) = \abs{\inner{\bm{f}(\Polygram)}{
\bm{f}(\Polygram')}}^{2}$, the per-coordinate restriction is a degree-$4$ trig
polynomial. Either way, it can be optimised in closed form (solve a quartic)
in $\Theta(1)$ time.
\end{remark}

\begin{theorem}[Coordinate-descent convergence]
\label{thm:cd-converge}
Block-coordinate descent over the $n$ Polygram factors, with each block solved
exactly, converges monotonically to a stationary point of $F$. The cost per
sweep is $\Theta(n N)$ for an objective involving $N$ reference Polygrams.
\end{theorem}
\begin{proof}
Monotone convergence follows from the standard theory of block-coordinate
descent on smooth functions: each block update strictly decreases $F$ unless
the block is already optimal, and the iterates lie in a compact parameter
space so admit a convergent subsequence. The cost claim follows from
Corollary~\ref{cor:complexity}: each per-block update costs $\Theta(N)$
(evaluate $N$ leave-one-out partial products) and there are $n$ blocks per
sweep. \qed
\end{proof}

% ---------------------------------------------------------------
\subsection{Riemannian gradient descent}
% ---------------------------------------------------------------

An alternative is to view the parameter space as the Riemannian manifold
$(\Sphere^{2})^{n}$ (Proposition~\ref{prop:riemannian}) and descend along
its geodesics.

\begin{proposition}[Geodesic step on $\Sphere^{2}$]
For a single factor $\bm{v}_{i} \in \Sphere^{2}$ with tangent vector
$\bm{t}_{i} \in T_{\bm{v}_{i}} \Sphere^{2}$ of norm $\norm{\bm{t}_{i}} = \tau$,
the geodesic step of length $\tau$ is
\[
\bm{v}'_{i} \;=\; \cos\tau\, \bm{v}_{i} + \sin\tau\, (\bm{t}_{i} / \tau).
\]
\end{proposition}

\begin{theorem}[Riemannian gradient descent]
\label{thm:riem-gd}
Riemannian gradient descent on $\Polygram$ with step size $\eta$ converges
to a stationary point of $F$ at rate $O(1/T)$ in the squared gradient norm,
where $T$ is the number of iterations, provided $\eta$ is small enough to
satisfy the standard descent lemma. Each iteration costs $\Theta(n N)$.
\end{theorem}

The proof is the standard convergence analysis for Riemannian gradient
methods on compact manifolds; nothing in it is Polygram-specific beyond the
per-iteration cost, which is given by Corollary~\ref{cor:total-gradient}.

% ---------------------------------------------------------------
\subsection{Greedy fitting / dictionary construction}
% ---------------------------------------------------------------

Given a target vector $\bm{u} \in \Sphere(\C^{D})$ stored only via an oracle
that returns $\inner{\bm{u}}{\bm{f}(\Polygram)}$ for any queried $\Polygram$,
we may wish to find a $\Polygram^{*}$ maximising this overlap. The natural
algorithm is greedy coordinate ascent.

\begin{center}
\setlength{\fboxsep}{8pt}
\fbox{%
\begin{minipage}{0.92\columnwidth}
\textbf{Algorithm 1 (Greedy Polygram fit).}\\[1mm]
\textbf{Input:} oracle access to $\bm{u}$; order $n$; tolerance $\epsilon$.\\
\textbf{Output:} $\Polygram^{*}$ maximising $\abs{\inner{\bm{u}}{
\bm{f}(\Polygram^{*})}}$.\\[1mm]
1.\; Initialise $\Polygram \leftarrow$ random Polygram.\\
2.\; \textbf{Repeat} until $\abs{F(\Polygram)}$ improves by less than
$\epsilon$:\\
3.\;\;\; \textbf{For each} factor $i = 1, \dots, n$:\\
4.\;\;\;\;\;\;\; Query the oracle with $n + 1$ probes that fix all factors
except $i$ at the current values; recover the marginal of $\bm{u}$ along
factor $i$ in closed form (Step~2 of Theorem~\ref{thm:identifiability}).\\
5.\;\;\;\;\;\;\; Update $\bm{v}_{i}$ to the optimal direction.\\
6.\;\textbf{Return} $\Polygram$.
\end{minipage}}
\end{center}

\begin{theorem}[Convergence of greedy fit]
\label{thm:greedy-fit}
Algorithm 1 produces a non-decreasing sequence of overlaps and converges to a
local maximum of $F$ at rate $O(1/T)$ in the squared overlap.
\end{theorem}

The number of oracle queries per outer iteration is $\Theta(n)$, and the
algorithm typically converges in $O(\log(1/\epsilon))$ outer iterations,
giving total query cost $O(n \log(1/\epsilon))$ in practice.

% ---------------------------------------------------------------
\subsection{Polygram merging}
% ---------------------------------------------------------------

In dictionary-maintenance applications, two Polygrams $\Polygram, \Polygram'$
with sufficiently high overlap may be \emph{merged} into a single Polygram
that approximates both. The natural choice is the geodesic midpoint in
$(\Sphere^{2})^{n}$.

\begin{proposition}[Merge fidelity]
\label{prop:merge}
Let $\Polygram, \Polygram'$ be Polygrams with $\rho_{i} =
\inner{\bm{v}_{i}}{\bm{w}_{i}}$ for each factor and overall overlap
$\rho = \prod_{i} \rho_{i}$. The geodesic midpoint $\Polygram^{*}$ in
$(\Sphere^{2})^{n}$ satisfies
\[
\inner{\bm{f}(\Polygram)}{\bm{f}(\Polygram^{*})}
   \;=\; \inner{\bm{f}(\Polygram^{*})}{\bm{f}(\Polygram')}
   \;=\; \prod_{i} \cos(\Theta_{i}/2)
\]
where $\Theta_{i}$ is the great-circle distance on $\Sphere^{2}$ between
$\bm{v}_{i}$ and $\bm{w}_{i}$. In particular,
\[
\inner{\bm{f}(\Polygram)}{\bm{f}(\Polygram^{*})} \;=\;
   \rho^{1/2}
\]
when all per-factor overlaps are positive reals.
\end{proposition}
\begin{proof}
The geodesic midpoint on $\Sphere^{2}$ between unit vectors $\bm{v}, \bm{w}$
at angular distance $\Theta$ has inner product $\cos(\Theta/2)$ with each
endpoint. Theorem~\ref{thm:fip} then factors the result. \qed
\end{proof}

The merge proposition shows that merging two Polygrams with overlap $\rho$
preserves overlap $\rho^{1/2}$ with both originals, a fact useful in
quantitatively designing merge-based dictionary maintenance schemes.

% ===============================================================
\section{Worked examples}
\label{sec:examples}
% ===============================================================

% ---------------------------------------------------------------
\subsection{Order 1 (qubits)}
% ---------------------------------------------------------------

For $n = 1$, a Polygram is just a single unit vector in $\C^{2}$. The Polygram
variety is $\Sphere(\C^{2})$ in its entirety (real dimension $3$), modulo
global phase it is $\C P^{1} \cong \Sphere^{2}$ (real dimension $2$).
Approximation is exact: every unit vector in $\C^{2}$ is a Polygram. The
factorised inner-product formula reduces to the Bloch inner product of
Lemma~\ref{lem:bloch-ip}, and the entire theory becomes the standard
geometry of the qubit Bloch ball.

% ---------------------------------------------------------------
\subsection{Order 2 (the simplest non-trivial case)}
% ---------------------------------------------------------------

For $n = 2$, the ambient sphere is $\Sphere(\C^{4})$ of real dimension $7$,
and the Polygram variety $\Sigma_{2}$ has real dimension $5$ (codimension
$2$).

The famous \emph{singlet state}
\[
\bm{u}_{\text{singlet}}
   = \frac{1}{\sqrt{2}}\big(\ket{01} - \ket{10}\big)
   = \frac{1}{\sqrt{2}}\,(0, 1, -1, 0)^{T}
\]
is a unit vector in $\C^{4}$ that is \emph{maximally non-separable}. One can
verify directly that the best Polygram fidelity for the singlet is
\[
\Lambda_{2}(\bm{u}_{\text{singlet}}) \;=\; \frac{1}{\sqrt{2}}.
\]
This is the smallest possible value of $\Lambda_{2}$ on $\Sphere(\C^{4})$,
and it matches the rate predicted by Theorem~\ref{thm:approx-lower} up to
constants: $n^{1/2} 2^{-n/2}$ at $n = 2$ gives $\sqrt{2}/2 = 1/\sqrt{2}$.
The example confirms that the worst-case bound of
Theorem~\ref{thm:approx-lower} is tight at small $n$.

In contrast, a product state such as $\ket{+}\ket{0} = \tfrac{1}{\sqrt 2}
(1, 0, 1, 0)^{T}$ has Polygram fidelity exactly $1$ because it is a Polygram:
take $\bm{v}_{1} = \bm{v}(\pi/4, 0)$ and $\bm{v}_{2} = \bm{v}(0, 0)$.

% ---------------------------------------------------------------
\subsection{Order 10 (representative of practical use)}
% ---------------------------------------------------------------

For $n = 10$ the ambient dimension is $D = 1024$. The Polygram variety has
real dimension $21$, an extraordinary compression: $21$ real numbers
parameterise an object that, viewed naively, requires $1024$ complex
coordinates ($2048$ real numbers). The compression ratio is $\sim 100\times$,
and inner products are roughly $50\times$ faster than the brute-force
implementation.

The worst-case approximation bound gives $\Lambda^{\min}_{10} \leq C \sqrt{10}
\cdot 2^{-5} = C \cdot 0.099$, i.e., there exists a unit vector in $\C^{1024}$
with best Polygram squared overlap below $\sim 0.01$. Most random vectors
in $\C^{1024}$ are this hard for Polygrams. Practical use therefore relies on
the target vectors being structured (typically themselves nearly Polygrams,
or low-rank mixtures of Polygrams).

% ===============================================================
\section{Open problems and future work}
\label{sec:open}
% ===============================================================

We close with a list of theoretical problems that would benefit from further
study. Each is at least as well-defined as the results above and most are
tractable with standard tools.

\begin{enumerate}[leftmargin=*, itemsep=4pt]
\item \textbf{Tight worst-case covering constant.} The bound of
      Theorem~\ref{thm:approx-lower} has the form $C n^{1/2} 2^{-n/2}$.
      A volume-comparison heuristic suggests the true constant $C$ can be
      taken close to $\sqrt{e/2}$; verifying this with explicit volumes of
      the Segre variety would sharpen the bound.

\item \textbf{Hardness of computing $\Lambda_{n}(\bm{u})$.} The geometric
      measure of entanglement is known to be NP-hard to compute for
      general inputs, but the worst case is irrelevant for parameter
      regimes where the input is itself a low-rank tensor. Establishing a
      polynomial-time algorithm or hardness result for ``$\bm{u}$ is given as
      a rank-$r$ MPS, output its Polygram fidelity to $1/\mathrm{poly}$'' is
      open.

\item \textbf{Sample complexity for recovery from noisy overlaps.} The
      construction of Theorem~\ref{thm:identifiability} uses $4n + 1$
      probes; the conjectured lower bound is $\Theta(n)$ for any
      $\mathrm{poly}(n, 1/\epsilon)$-accurate algorithm. A matching lower bound
      against an oblivious adversary is unknown.

\item \textbf{Polygram lattices.} Define a \emph{Polygram lattice} as the
      set of Polygrams whose Bloch angles lie on a regular grid in
      $[0, \pi/2] \times [0, 2\pi)$. How large must the grid be to
      $\epsilon$-cover the Polygram variety? The naive count is
      $\Theta(\epsilon^{-2n})$, but tighter packings via $A_{2}$- or
      $E_{8}$-style lattices in product spaces may help.

\item \textbf{Stability of dictionary recovery under cross-Polygram
      coherence.} Theorem~\ref{thm:noise-robust} bounds the noise
      sensitivity of a single Polygram's recovery; the corresponding
      multi-Polygram dictionary recovery (with $N$ unknowns and pairwise
      coherence $\mu$) is not yet quantified.

\item \textbf{Sharper MPS comparisons.} Table~\ref{tab:tradeoff} gives
      asymptotic costs but not constants. Empirical and analytic studies
      of the break-even point between $r = 1$ and $r = 2$ MPS on classes
      of structured tensors (e.g.\ low-Schmidt-rank random states) would
      indicate when Polygrams should be preferred.

\item \textbf{Generalised base alphabets.} Our definition fixes the per-site
      Hilbert space at $\C^{2}$. The same theory extends to $\C^{d}$ for
      any $d$, with $\Sigma_{n}$ a subvariety of $\C^{d^{n}}$ of real
      dimension $2(d-1) n$. The basic algebraic results carry over; the
      geometric and approximation bounds need to be re-derived with $d$ as
      a parameter.

\item \textbf{Polygrams over $\R$.} Real Polygrams (with $\bm{v}_{i} \in
      \Sphere^{1}$) are a strictly smaller family with $n$ real parameters
      and similar but simpler theory. The trade-off between real and
      complex Polygrams in applications that don't natively use complex
      numbers is open.

\item \textbf{Mixture-of-Polygrams expressivity.} The mixture family of
      Section~\ref{sec:mps} interpolates between Polygrams and full MPS.
      Characterising precisely which vectors admit a sparse (small-$r$)
      Polygram mixture but no small-bond MPS would settle the
      practical question of which compression family to choose.

\item \textbf{Polygram-friendly linear operators.} Local operators of the
      form $A = A_{1} \otimes \cdots \otimes A_{n}$ preserve Polygrams.
      Characterising the maximal set of linear operators on $\C^{D}$ that
      preserve Polygrams (or that preserve the family up to bounded
      Polygram fidelity loss) would identify the natural symmetry group
      of the data structure.
\end{enumerate}

% ===============================================================
\section*{Acknowledgements}
% ===============================================================

This manuscript is a theoretical treatment of an engineering primitive whose
roots lie in many decades of work on tensor decompositions, separable states,
and tensor networks. We acknowledge the broader literature on matrix product
states (Affleck, Kennedy, Lieb, Tasaki; Verstraete, Cirac), tensor trains
(Oseledets, Tyrtyshnikov), and the geometric measure of entanglement
(Wei, Goldbart) as the proper home for the results above. The Polygram
abstraction collects standard facts into a single self-contained object
intended for algorithmic use.

% ===============================================================
\begin{thebibliography}{99}
% ===============================================================

\bibitem{ostlund1995}
S. \"Ostlund and S. Rommer,
\emph{Thermodynamic limit of density matrix renormalization},
Phys.\ Rev.\ Lett.\ \textbf{75} (1995), 3537--3540.

\bibitem{verstraete2004}
F. Verstraete, D. Porras, and J. I. Cirac,
\emph{Density matrix renormalization group and periodic boundary conditions},
Phys.\ Rev.\ Lett.\ \textbf{93} (2004), 227205.

\bibitem{oseledets2011}
I. V. Oseledets,
\emph{Tensor-train decomposition},
SIAM J.\ Sci.\ Comput.\ \textbf{33} (2011), 2295--2317.

\bibitem{schollwock2011}
U. Schollw\"ock,
\emph{The density-matrix renormalization group in the age of matrix product
states},
Ann.\ Phys.\ (NY) \textbf{326} (2011), 96--192.

\bibitem{wei2003}
T.-C. Wei and P. M. Goldbart,
\emph{Geometric measure of entanglement and applications to bipartite and
multipartite quantum states},
Phys.\ Rev.\ A \textbf{68} (2003), 042307.

\bibitem{landsberg2012}
J. M. Landsberg,
\emph{Tensors: Geometry and Applications},
American Mathematical Society, 2012.

\bibitem{hackbusch2012}
W. Hackbusch,
\emph{Tensor Spaces and Numerical Tensor Calculus},
Springer, 2012.

\bibitem{evenbly2015}
G. Evenbly and G. Vidal,
\emph{Tensor network states and geometry},
J.\ Stat.\ Phys.\ \textbf{145} (2011), 891--918.

\bibitem{vershynin2018}
R. Vershynin,
\emph{High-Dimensional Probability},
Cambridge University Press, 2018.

\bibitem{hayden2006}
P. Hayden, D. Leung, and A. Winter,
\emph{Aspects of generic entanglement},
Comm.\ Math.\ Phys.\ \textbf{265} (2006), 95--117.

\bibitem{gross2006}
D. Gross, S. T. Flammia, and J. Eisert,
\emph{Most quantum states are too entangled to be useful as computational
resources},
Phys.\ Rev.\ Lett.\ \textbf{102} (2009), 190501.

\bibitem{absil2008}
P.-A. Absil, R. Mahony, and R. Sepulchre,
\emph{Optimization Algorithms on Matrix Manifolds},
Princeton University Press, 2008.

\bibitem{kolda2009}
T. G. Kolda and B. W. Bader,
\emph{Tensor decompositions and applications},
SIAM Review \textbf{51} (2009), 455--500.

\end{thebibliography}

\end{document}
